\exercice{Danger des infrarouges (20 minutes)}\points{8}

\begin{doc}{Les différents domaine d'infrarouge IR}
  Les émissions infrarouges sont classées en 3 domaines : infrarouge proche (A), infrarouge moyen (B) et infrarouge lointain (C).
  Ces domaines ont été établis à partir des propriétés d'absorption des tissus du corps humain.

  \begin{tableau}{|c |c |c |c |}
    Domaine & IR-A & IR-B & IR-C \\
    %
    Longueur d'onde &
    \qty{800}{\nm} $< \lambda <$ \qty{1400}{\nm} & 
    \qty{1400}{\nm} $< \lambda <$ \qty{3000}{\nm} & 
    \qty{3000}{\nm} $< \lambda <$ \qty{e6}{\nm} \\
    %
  \end{tableau}  
  
  Les \important{infrarouge-A (IR-A)} endommagent surtout la rétine, les yeux étant transparent pour ces longueurs d'ondes. La peau peut être brûlée superficiellement suite à une exposition prolongée au IR-A.
  Le cristallin peut être endommagé par des rayonnements \important{IR-A et IR-B} intenses.
  La cornée peut être endommagée par des rayonnements \important{IR-B et IR-C} très intenses et prolongés.
\end{doc}

\begin{doc}{Les métiers à risque IR}[\label{doc:metier_IR}]
  \begin{tableau}{|c|c|c|c|}
    Métier    & Source & Température en \unit{\degreeCelsius} & Durée d'exposition \\
    Boulanger & Four  & 250 & 5 h par jour \\
    \variationSujet{Soudeur}{Sidérurgiste}   & \variationSujet{Chalumeau}{Métal en fusion} & \variationSujet{3100}{2400} & 9 h par jour \\
    \variationSujet{Pompier}{Souffleur de verre}   & \variationSujet{Incendie}{Verre fondu} & \variationSujet{< 1500}{1130} & \variationSujet{10 h par semaine}{8 h par jour}
  \end{tableau}
\end{doc}

\begin{doc}{Loi de Wien et données}
  \begin{equation*}
    \lambda_{max} = \dfrac{\qty{2.9e-3}{\kelvin\m}}{T (\unit{\kelvin})}
  \end{equation*}

  \begin{listePoints}[2]
    \item \qty{1e-9}{\m} = \qty{1}{\nm}
  \end{listePoints}
\end{doc}

\question{
  Convertir les trois températures du document~\ref{doc:metier_IR} en kelvin.\points{1,5}
}{
  \begin{align*}
    \item  250 + \qty{273}{\degreeCelsius} = \qty{525}{\kelvin}
    \item \variationSujet
    {3100 + \qty{273}{\degreeCelsius} = \qty{3373}{\kelvin}}
    {2400 + \qty{273}{\degreeCelsius} = \qty{2673}{\kelvin}}
    \item \variationSujet
    {1500 + \qty{273}{\degreeCelsius} = \qty{1773}{\kelvin}}
    {1130 + \qty{273}{\degreeCelsius} = \qty{1403}{\kelvin}}
  \end{align*}
}

\question{
  Calculer la longueur d'onde maximale en utilisant le doc. 3 pour ces températures.\points{3}
}{
  \begin{align*}
    \item \lambda_\text{max}
    = \dfrac{\qty{2.9e-3}{\kelvin\m}}{\qty{525}{\kelvin}}
    = \qty{5.52e-6}{\m} = \qty{5520}{\nm}
    \item \lambda_\text{max}
    \variationSujet
    {= \dfrac{\qty{2.9e-3}{\kelvin\m}}{\qty{3373}{\kelvin}}
     = \qty{8.60e-7}{\m} = \qty{860}{\nm}}
    {= \dfrac{\qty{2.9e-3}{\kelvin\m}}{\qty{2673}{\kelvin}}
     = \qty{1.08e-6}{\m} = \qty{1080}{\nm}}
    \item \lambda_\text{max}
    \variationSujet
    {= \dfrac{\qty{2.9e-3}{\kelvin\m}}{\qty{1773}{\kelvin}}
     = \qty{1.63e-6}{\m} = \qty{1630}{\nm}}
    {= \dfrac{\qty{2.9e-3}{\kelvin\m}}{\qty{1403}{\kelvin}}
     = \qty{2.07e-6}{\m} = \qty{2070}{\nm}}
  \end{align*}
}

\question{
  Déduire pour chacun des métiers à quel type de rayonnement infrarouge ils sont exposés.\points{1,5}
}{
  IR-C pour le boulanger, IR-A pour le soudeur, IR-B pour le pompier.
}

\question{
  Conclure sur les risques liés à ces trois métiers et sur le type de protections à porter pour chaque métier.\points{2}
}{}